
\renewcommand{\Folder}{2025_SkinReplacement}
\renewcommand{\EntradaBibtex}{2025_SkinReplacement}

\begin{frame}{\citetitle{\EntradaBibtex} \footnotemark[1] (1)}

\begin{columns}
	\begin{column}{0.58\textwidth}
		\begin{itemize}
			\item Se desarrolló una aplicación móvil para Android que utiliza librerías de visión por computadora (OpenCV) para modificar el color de la piel en tiempo real.
			\item La aplicación detecta las áreas de piel mediante la combinación de los espacios de color HSV y YCrCb, permitiendo su aclarado u oscurecimiento con un control deslizante, además de alternar entre cámaras y capturar fotos con el efecto aplicado.
		\end{itemize}
	\end{column}
	\begin{column}{0.2\textwidth}
		 \begin{center}
		     \includegraphics[width=0.98\textwidth]{\Folder/figs/ejemplo_efecto_aplicado}
		 \end{center}
	\end{column}
	\begin{column}{0.2\textwidth}
		 \begin{center}
		     \includegraphics[width=0.98\textwidth]{\Folder/figs/ejemplo_mascara_piel}
		 \end{center}

	%\end{column}
	%\begin{column}{0.22\textwidth}  
		%\begin{center}
		 %\includegraphics[width=0.98\textwidth]{\Folder/figs/cuadroG}
		 %\end{center}
	\end{column}
\end{columns}
%\end{block} 
\footnotetext[1]{\fullcite{\EntradaBibtex}}
\end{frame}


\begin{frame}{\citetitle{\EntradaBibtex} \footnotemark[1] (2)}
%\begin{block}{Brazo Robótico \footnotemark} 
\begin{columns}
	\begin{column}{0.68\textwidth}
		\begin{itemize}
			\item Se logró implementar con éxito un algoritmo de detección de piel robusto, combinando los espacios de color HSV y YCrCb, con algunas limitaciones.  
            \item Las pruebas en tiempo real mostraron que el método es efectivo para identificar las áreas de piel en los fotogramas de la cámara, incluso bajo diferentes condiciones de iluminación. 
            \item Esto permitió generar una máscara binaria precisa que sirve como base para la aplicación del efecto.
		\end{itemize}
	\end{column}
	\begin{column}{0.3\textwidth}
		 \begin{center}
		     \includegraphics[width=0.98\textwidth]{\Folder/figs/dos_sujetos}
		 \end{center}
	\end{column}
	\begin{column}{0.02\textwidth}
		 \begin{center}
%		     \includegraphics[width=0.98\textwidth]{\Folder/figs/shownMunicipalityDetails}
		 \end{center}

	%\end{column}
	%\begin{column}{0.22\textwidth}  
		%\begin{center}
		 %\includegraphics[width=0.98\textwidth]{\Folder/figs/cuadroG}
		 %\end{center}
	\end{column}
\end{columns}
%\end{block} 
%\footnotetext[1]{\fullcite{\EntradaBibtex}}
\end{frame}

